\chapter{Cơ sở lý thuyết}

\section{Kiến trúc tập lệnh RISC-V}

\subsection{Lịch sử hình thành}

RISC-V là một kiến trúc tập lệnh (Instruction Set Architecture - ISA) 
dựa trên nguyên lý RISC (Reduced Instruction Set Computer), 
được đưa ra vào năm 2010 tại Phòng thí nghiệm Tính toán song song 
của Đại học California, Berkeley. Dự án được dẫn dắt bởi 
Giáo sư Krste Asanović cùng với sự tham gia của hai sinh viên cao học 
là Yunsup Lee và Andrew Waterman, dưới sự cố vấn của David Patterson – cha đẻ 
của kiến trúc RISC nguyên thủy.

Khác với các thế hệ RISC trước đó (như RISC-I, II, III, IV) vốn được 
thiết kế chủ yếu cho mục đích nghiên cứu học thuật, RISC-V ngay từ đầu đã 
được định hướng để trở thành một chuẩn ISA tiêu chuẩn cho cả nghiên cứu lẫn 
thương mại hóa. Mục tiêu của nhóm thiết kế là tạo ra một tập lệnh "sạch", 
không chịu gánh nặng của các vấn đề tương thích ngược từ các kiến trúc cũ kỹ 
tồn tại hàng thập kỷ, đồng thời hoàn toàn mở và miễn phí.

Năm 2015, RISC-V Foundation (nay là RISC-V International) được thành 
lập để quản lý chuẩn hóa và thúc đẩy hệ sinh thái phần cứng/phần mềm. 
Sự ra đời của RISC-V đã phá vỡ thế độc quyền lâu năm của các kiến trúc 
đóng như x86 (Intel) và ARM, mở ra kỷ nguyên "phần cứng mã nguồn mở" 
(Open Source Hardware).

\subsection{Các đặc điểm nổi bật}
So với các kiến trúc thương mại phổ biến, RISC-V sở hữu các đặc điểm kỹ 
thuật và pháp lý vượt trội:
\begin{itemize}
    \item \textbf{Tính mở và Miễn phí (Open \& Free):} RISC-V được phát hành dưới giấy phép BSD. Bất kỳ cá nhân hay tổ chức nào cũng có thể thiết kế, sản xuất và thương mại hóa chip RISC-V mà không phải trả phí bản quyền (Royalty-free). Điều này giúp giảm đáng kể chi phí khởi tạo cho các dự án khởi nghiệp và nghiên cứu.
    \item \textbf{Tính Mô đun:} Thay vì bắt buộc phải hiện thực toàn bộ tập lệnh khổng lồ, RISC-V có lõi là tập lệnh số nguyên cơ sở (Base Integer ISA) rất nhỏ gọn. Các tính năng nâng cao như số thực, nhân chia, hay nén lệnh (Compressed) được coi là các phần mở rộng tùy chọn. Điều này cho phép tối ưu hóa diện tích vi mạch cho các ứng dụng nhúng nhỏ gọn.
    \item \textbf{Khả năng mở rộng (Extensibility):} RISC-V dành riêng một phần không gian mã hóa lệnh (Opcode space) cho người dùng tự định nghĩa. Điều này đặc biệt quan trọng đối với các ứng dụng Edge AI, nơi các kỹ sư có thể thêm các lệnh tính toán ma trận hoặc tích chập chuyên biệt để tăng tốc phần cứng mà không vi phạm chuẩn ISA.
    \item \textbf{Kiến trúc Load/Store thuần túy:} Tương tự các kiến trúc RISC khác, RISC-V tách biệt hoàn toàn việc truy cập bộ nhớ và xử lý dữ liệu, giúp đơn giản hóa thiết kế phần cứng của Pipeline.
\end{itemize}

\subsection{Cấu trúc tập lệnh cơ sở và các phần mở rộng}

\subsubsection{Tập lệnh số nguyên cơ sở (Base Integer ISA)}
Nền tảng của RISC-V là tập lệnh số nguyên cơ sở, được ký hiệu là \textbf{RV32I} (cho bản 32-bit), RV64I (64-bit) hoặc RV128I (128-bit). Trong phạm vi đồ án này, nhóm tập trung vào \textbf{RV32I}, bao gồm 47 lệnh cơ bản.

\textbf{Hệ thống thanh ghi (Register File):}
RV32I cung cấp 32 thanh ghi đa năng (General Purpose Registers - GPR), mỗi thanh ghi rộng 32-bit, được đánh số từ $x0$ đến $x31$. Ngoài ra còn có một thanh ghi đếm chương trình (Program Counter - PC).

Theo quy ước gọi hàm (Calling Convention) của RISC-V ABI, các thanh ghi này có vai trò cụ thể như mô tả trong Bảng \ref{tab:riscv_registers}:

\begin{table}[h]
\centering
\begin{tabular}{|l|l|l|l|}
\hline
\textbf{Tên} & \textbf{ABI Name} & \textbf{Mô tả} & \textbf{Lưu bởi} \\ \hline
x0 & zero & Luôn có giá trị bằng 0 (Hardwired Zero) & - \\ \hline
x1 & ra & Địa chỉ trả về (Return Address) & Caller \\ \hline
x2 & sp & Con trỏ ngăn xếp (Stack Pointer) & Callee \\ \hline
x3 & gp & Con trỏ toàn cục (Global Pointer) & - \\ \hline
x4 & tp & Con trỏ luồng (Thread Pointer) & - \\ \hline
x5-x7 & t0-t2 & Thanh ghi tạm thời (Temporaries) & Caller \\ \hline
x8 & s0/fp & Thanh ghi lưu trữ/Con trỏ khung (Frame Ptr) & Callee \\ \hline
x9 & s1 & Thanh ghi lưu trữ (Saved Register) & Callee \\ \hline
x10-x11 & a0-a1 & Tham số hàm / Giá trị trả về & Caller \\ \hline
x12-x17 & a2-a7 & Tham số hàm (Function Arguments) & Caller \\ \hline
x18-x27 & s2-s11 & Thanh ghi lưu trữ (Saved Registers) & Callee \\ \hline
x28-x31 & t3-t6 & Thanh ghi tạm thời (Temporaries) & Caller \\ \hline
\end{tabular}
\caption{Quy ước sử dụng thanh ghi trong RISC-V (ABI)}
\label{tab:riscv_registers}
\end{table}

\subsubsection{Định dạng lệnh (Instruction Formats)}
Để đơn giản hóa bộ giải mã, các lệnh trong RV32I có độ dài cố định 32-bit và được căn chỉnh bộ nhớ (aligned) theo địa chỉ chia hết cho 4. RISC-V định nghĩa 6 định dạng lệnh cơ bản:

\begin{itemize}
    \item \textbf{R-type (Register):} Dùng cho các phép toán giữa các thanh ghi (VD: cộng, trừ, logic). Cấu trúc gồm: opcode, rd (đích), funct3, rs1 (nguồn 1), rs2 (nguồn 2), funct7.
    \item \textbf{I-type (Immediate):} Dùng cho các phép toán với hằng số ngắn hoặc lệnh Load. Gồm: opcode, rd, funct3, rs1, imm[11:0].
    \item \textbf{S-type (Store):} Dùng cho lệnh lưu dữ liệu ra bộ nhớ (Store). Hằng số immediate được chia làm 2 phần để giữ vị trí các thanh ghi nguồn cố định.
    \item \textbf{B-type (Branch):} Dùng cho các lệnh rẽ nhánh có điều kiện. Tương tự S-type nhưng immediate được mã hóa khác để mở rộng phạm vi nhảy.
    \item \textbf{U-type (Upper Immediate):} Dùng cho lệnh thao tác với 20-bit cao của thanh ghi (VD: LUI, AUIPC).
    \item \textbf{J-type (Jump):} Dùng cho lệnh nhảy không điều kiện (JAL).
\end{itemize}


\subsubsection{Các chế độ đặc quyền}
RISC-V định nghĩa ba chế độ đặc quyền để bảo vệ hệ thống:
\begin{enumerate}
    \item \textbf{Machine Mode (M-Mode):} Chế độ có quyền cao nhất, truy cập toàn bộ tài nguyên phần cứng. Các vi điều khiển đơn giản (bare-metal) chỉ cần chạy ở chế độ này.
    \item \textbf{Supervisor Mode (S-Mode):} Dùng cho hệ điều hành (như Linux, FreeRTOS) để quản lý bộ nhớ ảo và bảo vệ tài nguyên.
    \item \textbf{User Mode (U-Mode):} Chế độ có quyền thấp nhất, dành cho các ứng dụng người dùng.
\end{enumerate}
Trong đồ án này, lõi RISC-V chủ yếu sẽ hoạt động ở M-Mode để tương tác trực tiếp với phần cứng FPGA và bộ tăng tốc.


\subsection{So sánh kiến trúc RISC và CISC}

Để làm rõ ưu thế của việc lựa chọn RISC-V cho đề tài, cần đặt nó vào hệ quy chiếu so sánh với kiến trúc CISC (Complex Instruction Set Computer) – triết lý thiết kế đối lập vốn thống trị thị trường máy tính cá nhân (PC) nhiều thập kỷ qua.

\subsubsection{Triết lý thiết kế}
Sự khác biệt căn bản giữa RISC và CISC nằm ở cách phân chia khối lượng công việc giữa phần cứng (Hardware) và phần mềm (Software/Compiler):

\begin{itemize}
    \item \textbf{CISC:} Tập trung vào việc giảm "khoảng cách ngữ nghĩa" (semantic gap) giữa ngôn ngữ bậc cao và mã máy. CISC cung cấp các lệnh phức tạp, đa năng (ví dụ: một lệnh duy nhất có thể thực hiện nạp dữ liệu, cộng và lưu kết quả). Điều này giúp giảm kích thước mã chương trình (Code size) nhưng làm tăng độ phức tạp của phần cứng giải mã lệnh.
    \item \textbf{RISC:} Chuyển gánh nặng tối ưu hóa sang trình biên dịch (Compiler). Tập lệnh chỉ bao gồm các chỉ thị đơn giản, thực thi nhanh. Các tác vụ phức tạp được thực hiện bằng cách kết hợp nhiều lệnh đơn giản. Điều này giúp phần cứng đơn giản hơn, tiết kiệm diện tích silicon để dành cho bộ nhớ đệm (Cache) hoặc các lõi xử lý song song.
\end{itemize}



\subsubsection{So sánh RISC và CISC}
Bảng \ref{tab:risc_vs_cisc} tóm tắt các đặc điểm kỹ thuật đối lập giữa hai kiến trúc \cite{risc-v}:

\begin{table}[h]
\centering
\begin{tabular}{|p{4cm}|p{5.5cm}|p{5.5cm}|}
\hline
\textbf{Tiêu chí} & \textbf{RISC (RISC-V, ARM)} & \textbf{CISC (x86)} \\ \hline
\textbf{Độ phức tạp lệnh} & Đơn giản, thực hiện tác vụ nhỏ. & Phức tạp, một lệnh làm nhiều việc (truy cập bộ nhớ, tính toán). \\ \hline
\textbf{Thời gian thực thi} & Hướng tới 1 chu kỳ máy/lệnh (CPI $\approx$ 1). & Đa chu kỳ (Multi-cycle), số chu kỳ thay đổi tùy lệnh. \\ \hline
\textbf{Định dạng lệnh} & Cố định (thường là 32-bit), giải mã nhanh. & Thay đổi (Variable length), giải mã phức tạp. \\ \hline
\textbf{Truy cập bộ nhớ} & Chỉ qua lệnh Load/Store. & Các lệnh tính toán có thể thao tác trực tiếp với bộ nhớ. \\ \hline
\textbf{Số lượng thanh ghi} & Nhiều (32 hoặc hơn) để giảm truy cập RAM. & Ít thanh ghi chuyên dụng hơn. \\ \hline
\textbf{Pipiline} & Dễ dàng hiện thực và đạt hiệu quả cao. & Khó hiện thực do độ dài lệnh không đồng nhất. \\ \hline
\textbf{Kích thước chương trình} & Lớn hơn (cần nhiều lệnh để làm một việc). & Nhỏ gọn hơn (mật độ mã cao). \\ \hline
\end{tabular}
\caption{So sánh đặc điểm kỹ thuật giữa RISC và CISC}
\label{tab:risc_vs_cisc}
\end{table}

\subsection{Các phần mở rộng tiêu chuẩn (Standard Extensions)}
Để tăng cường khả năng xử lý, RISC-V cung cấp các phần mở rộng được chuẩn hóa. Một vi xử lý hỗ trợ đầy đủ các mở rộng cơ bản thường được ký hiệu là \textbf{RV32G} (General), trong đó G = I + M + A + F + D.

\begin{itemize}
    \item \textbf{M (Multiplication/Division):} Bổ sung các lệnh nhân và chia số nguyên phần cứng. Nếu không có mở rộng này, CPU phải thực hiện nhân/chia bằng phần mềm rất chậm. Đây là mở rộng bắt buộc cho các ứng dụng AI.
    \item \textbf{A (Atomic):} Cung cấp các lệnh truy cập bộ nhớ nguyên tử (Read-Modify-Write), cần thiết để hỗ trợ hệ điều hành đa nhiệm (như Linux) và đồng bộ hóa đa luồng.
    \item \textbf{F (Single-Precision Floating Point):} Hỗ trợ tính toán số thực dấu chấm động 32-bit (chuẩn IEEE 754-2008), bổ sung thêm 32 thanh ghi số thực $f0-f31$.
    \item \textbf{D (Double-Precision Floating Point):} Hỗ trợ số thực 64-bit.
    \item \textbf{C (Compressed Instructions):} Hỗ trợ các lệnh nén độ dài 16-bit. Mở rộng này giúp giảm kích thước chương trình (code size) từ 25-30\%, cực kỳ hữu ích cho các thiết bị nhúng có bộ nhớ hạn chế.
    \item \textbf{V (Vector Operations):} Mở rộng xử lý vector, cho phép thực hiện một lệnh trên nhiều dữ liệu (SIMD). Đây là mở rộng quan trọng nhất đối với các ứng dụng học sâu (Deep Learning) và xử lý ảnh hiện đại, tuy nhiên trong phạm vi đồ án này, chúng ta sẽ thiết kế bộ tăng tốc riêng thay vì dùng V-extension tiêu chuẩn.
\end{itemize}

\subsection{Sự phù hợp của RISC đối với Edge AI}
Mặc dù CISC vẫn thống trị mảng PC/Server nhờ tính tương thích ngược, RISC (đặc biệt là RISC-V) lại chiếm ưu thế tuyệt đối trong lĩnh vực Hệ thống nhúng và IoT vì các lý do sau:
\begin{enumerate}
    \item \textbf{Hiệu quả năng lượng (Power Efficiency):} Do phần cứng giải mã đơn giản, các chip RISC tiêu thụ ít điện năng hơn đáng kể, phù hợp cho các thiết bị chạy pin.
    \item \textbf{Diện tích nhỏ:} Thiết kế RISC chiếm ít tài nguyên LUTs/Flip-flops hơn khi triển khai trên FPGA, để dành không gian cho các bộ tăng tốc AI.
    \item \textbf{Dễ dàng tùy biến:} Cấu trúc lệnh đơn giản của RISC giúp các nhà thiết kế dễ dàng thêm các lệnh tùy chỉnh (Custom Instructions) cho AI mà không làm phá vỡ kiến trúc đường ống, điều rất khó thực hiện trên CISC.
\end{enumerate}

%=================================================
%
%=================================================
\section{Kiến trúc Vi xử lý PULP (Pulpino/PULPissimo)}

Trong khi mục 2.1 trình bày về chuẩn kiến trúc tập lệnh (ISA) RISC-V trên lý thuyết, phần này sẽ đi sâu vào việc hiện thực hóa phần cứng cụ thể được sử dụng trong đồ án: nền tảng PULP (Parallel Ultra Low Power). Đây là bước chuyển quan trọng từ lý thuyết sang thực tiễn triển khai hệ thống SoC trên FPGA.

\subsection{Tổng quan về nền tảng PULP}
PULP là một dự án mã nguồn mở được khởi xướng bởi sự hợp tác giữa ETH Zurich (Thụy Sĩ) và Đại học Bologna (Ý). Mục tiêu tối thượng của PULP không chỉ là tạo ra một vi xử lý hoạt động được, mà là đạt được hiệu quả năng lượng cực cao (Ultra-low power) và hiệu năng tính toán vượt trội cho các tác vụ xử lý tín hiệu số (DSP) tại biên.

Pulpino (hoặc phiên bản nâng cấp là PULPissimo) là một hệ thống SoC đơn nhân (Single-core SoC) thuộc họ PULP. Kiến trúc tổng thể của nó bao gồm một lõi xử lý trung tâm, hệ thống bộ nhớ SRAM tốc độ cao, và một tập hợp các thiết bị ngoại vi phong phú (SPI, I2C, UART, GPIO) được kết nối thông qua hệ thống bus AXI4 và APB. Điểm đặc biệt là kiến trúc này loại bỏ bộ nhớ Cache truyền thống để thay thế bằng bộ nhớ Scratchpad (TCDM) nhằm đảm bảo tính xác định về thời gian (determinism) cho các ứng dụng thời gian thực.

\subsection{Nhân xử lý (Core Implementation): CV32E40P}
Trái tim của hệ thống là nhân vi xử lý \textbf{CV32E40P} (tên mã cũ là RI5CY). Đây là một hiện thực vi kiến trúc (Micro-architecture) 4 tầng đường ống (4-stage pipeline), được thiết kế tối ưu cho hiệu năng trên năng lượng.

\subsubsection{Cấu trúc đường ống lệnh}
Đường ống của CV32E40P bao gồm 4 giai đoạn:
\begin{enumerate}
    \item \textbf{Instruction Fetch (IF):} Nạp lệnh từ bộ nhớ lệnh hoặc bộ đệm lệnh (Instruction Cache/Buffer).
    \item \textbf{Instruction Decode (ID):} Giải mã lệnh, đọc các toán hạng từ tập thanh ghi (Register File). Tại đây, các lệnh mở rộng DSP cũng được giải mã.
    \item \textbf{Execute (EX):} Thực thi phép toán. Giai đoạn này chứa bộ ALU (Arithmetic Logic Unit), bộ nhân (Multiplier) và các đơn vị tính toán vector.
    \item \textbf{Write Back (WB):} Ghi kết quả ngược lại vào tập thanh ghi.
\end{enumerate}

Với độ sâu đường ống ngắn, CV32E40P giảm thiểu được chi phí phạt (penalty) khi dự đoán sai nhánh (Branch Misprediction), đồng thời giữ cho thiết kế phần cứng nhỏ gọn, phù hợp với tài nguyên hạn chế của FPGA.

\subsection{Các mở rộng tập lệnh DSP (DSP Extensions)}
Sức mạnh thực sự của PULP nằm ở bộ mở rộng tập lệnh (Custom Extensions) giúp tăng tốc độ xử lý các thuật toán AI lên nhiều lần so với chuẩn RISC-V thông thường.

\subsubsection{Hardware Loops (Vòng lặp phần cứng)}
Trong các vi xử lý thông thường, mỗi vòng lặp `for` hoặc `while` đều tốn các chu kỳ máy cho việc tăng biến đếm, so sánh điều kiện và nhảy về đầu vòng lặp. PULP loại bỏ hoàn toàn lãng phí này bằng cách sử dụng các thanh ghi phần cứng chuyên dụng để quản lý vòng lặp.

\subsubsection{SIMD và Post-increment}
\begin{itemize}
    \item \textbf{Packed-SIMD:} Cho phép thực hiện cộng/nhân song song trên các dữ liệu nhỏ (8-bit hoặc 16-bit) được gói trong thanh ghi 32-bit. Đây là tính năng then chốt cho lượng tử hóa (Quantization) trong AI.
    \item \textbf{Post-increment Load/Store:} Tự động cập nhật địa chỉ con trỏ ngay sau khi truy cập bộ nhớ, giảm số lượng lệnh `ADD` cần thiết.
\end{itemize}

\subsubsection{So sánh mã nguồn Assembly}
Để minh họa hiệu quả, ta xem xét đoạn mã tính tích vô hướng (Dot Product) - phép toán cơ bản nhất của lớp Tích chập (Convolution) trong CNN:

\begin{table}[h]
\centering
\begin{tabular}{|p{7cm}|p{7cm}|}
\hline
\textbf{RISC-V Chuẩn (RV32IM)} & \textbf{PULP Optimized (DSP Extension)} \\ \hline
\ttfamily
\small
loop: \\
\ \ lw \ \ \ t0, 0(a0) \ \ \ // Load a[i] \newline
\ \ lw \ \ \ t1, 0(a1) \ \ \ // Load b[i] \newline
\ \ mul \ \ t2, t0, t1 \ \ // t2 = a[i]*b[i] \newline
\ \ add \ \ a2, a2, t2 \ \ // sum += t2 \newline
\ \ addi \ a0, a0, 4 \ \ \ // Incr pointer a \newline
\ \ addi \ a1, a1, 4 \ \ \ // Incr pointer b \newline
\ \ addi \ a3, a3, -1 \ \ // Decr counter \newline
\ \ bne \ \ a3, zero, loop // Branch check
& 
\ttfamily
\small
\ \ // Setup Hardware Loop \newline
\ \ lp.setup x0, a3, end\_loop \newline
\ \ \newline
\ \ // Body (SIMD \& Post-inc) \newline
\ \ pv.sdotsp.h \ a2, a0, a1 \newline
\ \ // 1 lenh lam tat ca viec tren \newline
\ \ \newline
end\_loop: \newline
\ \ // Tu dong quay lai
\\ \hline
\textbf{Kết quả:} Tốn 7-8 chu kỳ/phần tử. & \textbf{Kết quả:} Tốn 1 chu kỳ/phần tử (với SIMD). \\ \hline
\end{tabular}
\caption{So sánh mã Assembly giữa RV32IM và PULP DSP cho Dot Product}
\end{table}

Qua ví dụ trên, có thể thấy kiến trúc PULP có thể tăng tốc độ tính toán 
lên gấp 4-8 lần so với kiến trúc chuẩn nhờ giảm số lượng lệnh (Instruction Count Reduction).

\subsection{Hệ thống bộ nhớ và Ngoại vi thông minh}

\subsubsection{Bộ nhớ TCDM (Tightly Coupled Data Memory)}
Thay vì sử dụng bộ nhớ Cache có độ trễ không đoán trước được, 
PULP sử dụng mô hình bộ nhớ TCDM. Đây là một vùng SRAM đa ngân hàng (Multi-banked SRAM).
\begin{itemize}
    \item \textbf{Cấu trúc Banked:} Bộ nhớ được chia thành nhiều ngân hàng độc lập (ví dụ: 8 hoặc 16 banks). Điều này cho phép lõi CPU và bộ điều khiển DMA truy cập đồng thời vào các vùng nhớ khác nhau mà không xảy ra xung đột (Memory Contention).
    \item \textbf{Độ trễ thấp:} Việc truy cập vào TCDM chỉ mất 1 chu kỳ máy (Single-cycle access), đảm bảo dữ liệu luôn sẵn sàng cho bộ tăng tốc AI.
\end{itemize}

\subsubsection{Bộ điều khiển uDMA (micro-DMA)}
Trong các hệ thống AI tại biên, việc di chuyển dữ liệu (từ Camera vào RAM, hoặc từ RAM ra bộ tăng tốc) tiêu tốn năng lượng hơn nhiều so với việc tính toán.
PULP giải quyết vấn đề này bằng uDMA - một bộ điều khiển truy cập bộ nhớ trực tiếp thông minh. uDMA có khả năng tự động chuyển luồng dữ liệu lớn từ các ngoại vi (như Camera qua giao tiếp CPI, hoặc cảm biến qua SPI) trực tiếp vào bộ nhớ TCDM mà hoàn toàn không cần CPU can thiệp. CPU có thể chuyển sang chế độ ngủ (Sleep mode) trong khi dữ liệu đang được nạp, giúp tiết kiệm năng lượng tối đa.

\subsection{Đơn vị quản lý sự kiện (Event Unit)}
Khác với các vi xử lý ARM sử dụng bộ điều khiển ngắt lồng nhau (NVIC) phức tạp, PULP sử dụng một Event Unit đơn giản hóa để quản lý năng lượng. Event Unit cho phép tắt hoàn toàn xung nhịp (Clock Gating) cấp cho nhân CPU khi hệ thống chờ sự kiện (ví dụ: chờ kết quả từ bộ tăng tốc AI), và đánh thức CPU ngay lập tức khi có sự kiện xảy ra. Cơ chế này cực kỳ quan trọng đối với thiết kế SoC chạy bằng pin.

%=================================================
%^^^^^^^^^^^^^^^^^^^^^^^^^^^^^^^^^^^^^^^^^^^^^^^^^^
%=================================================
\newpage
\section{Nền tảng phần cứng FPGA và System-on-Chip}

Để hiện thực hóa một hệ thống Edge AI có khả năng tùy biến cao, 
việc lựa chọn nền tảng phần cứng là yếu tố then chốt. 
Phần này trình bày từ các khái niệm cơ bản về công nghệ FPGA, 
kiến trúc hệ thống trên vi mạch (SoC) đến chi tiết kỹ thuật của 
kit phát triển Kria KV260 được sử dụng trong đồ án.

\subsection{Công nghệ FPGA (Field Programmable Gate Array)}
\subsubsection{Khái niệm và Kiến trúc cơ bản}
FPGA (Field Programmable Gate Array - Mảng cổng lập trình dạng trường) 
là một loại vi mạch bán dẫn đặc biệt cho phép người dùng tái cấu hình 
phần cứng sau khi đã sản xuất. Khác với ASIC (Application Specific Integrated Circuit) 
vốn có chức năng cố định vĩnh viễn, 
FPGA bao gồm một ma trận các khối logic khả trình (Configurable Logic Blocks - CLBs) 
được kết nối với nhau thông qua mạng lưới dây dẫn có thể lập trình 
(Programmable Interconnects).\cite{fpga1}

Thành phần cấu tạo cơ bản của FPGA bao gồm:
\begin{itemize}
    \item \textbf{Look-Up Table (LUT):} Đơn vị cơ bản nhất dùng để hiện thực các hàm logic tổ hợp (Boolean logic). Một LUT N-đầu vào có thể coi như một bộ nhớ RAM nhỏ lưu trữ bảng chân trị của hàm số.
    \item \textbf{Flip-Flops (FF):} Phần tử nhớ dùng để lưu trạng thái, kết hợp với LUT để tạo thành các mạch tuần tự (Sequential Logic).
    \item \textbf{DSP Slices (Digital Signal Processing):} Các khối phần cứng chuyên dụng thực hiện phép nhân và cộng (Multiply-Accumulate - MAC) tốc độ cao. Đây là tài nguyên quan trọng nhất đối với các ứng dụng xử lý tín hiệu số và Trí tuệ nhân tạo.
    \item \textbf{Block RAM (BRAM):} Các khối bộ nhớ tĩnh (SRAM) nằm rải rác trong chip, cung cấp bộ nhớ đệm cục bộ với băng thông cực lớn và độ trễ thấp.
\end{itemize}

\subsubsection{Ưu điểm trong ứng dụng AI}
Nhờ khả năng xử lý song song ở mức phần cứng, FPGA có thể thực hiện hàng nghìn phép tính nhân-cộng đồng thời, phù hợp với đặc thù tính toán ma trận dày đặc của mạng CNN. Hơn nữa, khả năng tùy biến luồng dữ liệu (Dataflow) giúp loại bỏ các nút thắt cổ chai về bộ nhớ thường gặp trên các kiến trúc Von Neumann truyền thống.

\subsection{System on Chip - SoC}
Trong các thiết kế hiện đại, FPGA hiếm khi đứng độc lập mà thường được tích hợp trong một Hệ thống trên vi mạch (SoC). Kiến trúc này, thường được gọi là SoC FPGA hoặc MPSoC (Multi-Processor SoC), kết hợp hai thế giới:

\begin{enumerate}
    \item \textbf{Hệ thống xử lý (Processing System - PS):} Bao gồm các nhân vi xử lý cứng (Hard-core) như ARM Cortex-A, chạy hệ điều hành (Linux/RTOS) để quản lý kết nối mạng, giao diện người dùng và các tác vụ cấp cao.
    \item \textbf{Phần logic lập trình được (Programmable Logic - PL):} Chính là phần FPGA, dùng để thực thi các tác vụ tính toán nặng, xử lý tín hiệu thời gian thực hoặc hiện thực các giao tiếp tùy chỉnh.
\end{enumerate}

Sự kết hợp này cho phép thực hiện đồng thiết kế phần cứng/phần mềm (Hardware/Software Co-design). Trong đồ án này, phần PL sẽ được sử dụng để nhúng lõi RISC-V mềm (Soft-core) và bộ tăng tốc AI, trong khi phần PS đóng vai trò giám sát và nạp chương trình.

\subsection{Nền tảng Kria KV260 Vision AI}
Để triển khai thực nghiệm, đồ án sử dụng bộ công cụ \textbf{Kria KV260 Vision AI Starter Kit} của hãng Xilinx. Đây là nền tảng phần cứng được tối ưu hóa đặc biệt cho các ứng dụng thị giác máy tính tại biên.

\subsubsection{Kiến trúc System on Module (SOM)}
Điểm khác biệt của dòng Kria so với các kit FPGA truyền thống là việc áp dụng kiến trúc SOM. Kit KV260 bao gồm hai phần tách biệt:
\begin{itemize}
    \item \textbf{K26 SOM:} Mô-đun tính toán trung tâm kích thước nhỏ gọn, chứa chip xử lý chính, bộ nhớ RAM (4GB DDR4), bộ nhớ Flash (16GB eMMC) và các mạch quản lý nguồn (PMIC).
    \item \textbf{Carrier Card:} Bo mạch đế cung cấp các giao tiếp ngoại vi như Ethernet, USB 3.0, HDMI, DisplayPort và các cổng kết nối Camera (MIPI CSI-2).
\end{itemize}

\subsubsection{Chip Zynq UltraScale+ MPSoC}
Trái tim của K26 SOM là chip \textbf{XCK26 Zynq UltraScale+ MPSoC}. Đây là dòng chip cao cấp được sản xuất trên tiến trình 16nm FinFET+, cung cấp tài nguyên dồi dào cho thiết kế:
\begin{itemize}
    \item \textbf{Processing System (PS):}
        \begin{itemize}
            \item Quad-core ARM Cortex-A53 (lên đến 1.5GHz) cho ứng dụng hệ thống.
            \item Dual-core ARM Cortex-R5F cho xử lý thời gian thực.
            \item GPU Mali-400 MP2 cho xử lý đồ họa.
        \end{itemize}
    \item \textbf{Programmable Logic (PL):}
        \begin{itemize}
            \item \textbf{Logic Cells:} 256,200 ô logic (tương đương khoảng 1.2 triệu cổng ASIC) - Đủ lớn để chứa lõi RISC-V phức tạp.
            \item \textbf{DSP Slices:} 1,248 khối DSP48E2 - Cung cấp năng lực tính toán lên tới hàng nghìn tỷ phép tính mỗi giây (TOPS) cho AI.
            \item \textbf{Block RAM:} 144 khối (36Kb mỗi khối) và 64 khối UltraRAM, cung cấp băng thông nội bộ cực lớn.
        \end{itemize}
\end{itemize}

\subsection{Hệ thống thu nhận hình ảnh}
Một thành phần không thể thiếu của hệ thống Edge AI thị giác máy tính là phân hệ thu nhận hình ảnh. KV260 hỗ trợ chuẩn giao tiếp công nghiệp \textbf{MIPI CSI-2}, cho phép kết nối với các cảm biến ảnh chất lượng cao.

\subsubsection{Luồng xử lý Video trên FPGA (Video Pipeline)}
Do dữ liệu ảnh thô không thể sử dụng ngay cho mạng nơ-ron, một chuỗi xử lý (Pipeline) phần cứng được thiết kế trên phần PL của FPGA:
\begin{enumerate}
    \item \textbf{MIPI CSI-2 RX Subsystem:} Khối IP cứng nhận tín hiệu vật lý từ camera, tách xung nhịp và giải mã gói tin dữ liệu.
    \item \textbf{Image Signal Processor (ISP):} Thực hiện các thuật toán tiền xử lý quan trọng:
        \begin{itemize}
            \item \textit{Demosaicing:} Nội suy màu để chuyển từ định dạng Bayer sang RGB.
            \item \textit{Auto White Balance / Auto Exposure:} Tự động cân bằng trắng và phơi sáng.
            \item \textit{Color Space Conversion:} Chuyển đổi không gian màu (ví dụ: RGB sang YUV) nếu cần thiết.
        \end{itemize}
    \item \textbf{Video DMA (VDMA):} Sử dụng giao thức AXI4-Stream để ghi luồng dữ liệu video đã xử lý trực tiếp vào bộ nhớ DDR4 thông qua bộ điều khiển bộ nhớ (Memory Controller).
\end{enumerate}

Tại thời điểm này, hình ảnh đã sẵn sàng trong bộ nhớ hệ thống. Core RISC-V sẽ truy cập vùng nhớ này thông qua bus AXI4 để lấy dữ liệu đầu vào cho quá trình suy luận AI.
%=================================================
%
%=================================================
\newpage
\section{Tổng quan về Trí tuệ Nhân tạo (AI)}
\label{sec:ai_overview}

\subsection{Các Mô hình AI (AI Models)}
Trong lĩnh vực thị giác máy tính và hệ thống nhúng, các mô hình AI thường được phân loại dựa trên chức năng và kiến trúc. Việc lựa chọn mô hình phù hợp là yếu tố quyết định đến hiệu năng của hệ thống Edge AI.

\subsubsection{Mô hình Phân loại (Classification)}
Đây là dạng mô hình cơ bản nhất, nhiệm vụ là gán nhãn cho một hình ảnh đầu vào. Ví dụ: xác định xem trong ảnh có người hay không. Các kiến trúc tiêu biểu bao gồm MobileNet, SqueezeNet, và ResNet. Đối với các thiết bị biên, MobileNet thường được ưu tiên nhờ kiến trúc Depthwise Separable Convolution giúp giảm đáng kể khối lượng tính toán.

\subsubsection{Mô hình Phát hiện đối tượng (Object Detection)}
Phức tạp hơn phân loại, mô hình này không chỉ xác định đối tượng là gì mà còn định vị vị trí của nó thông qua các khung bao (Bounding Box). Các đại diện nổi bật gồm dòng YOLO (You Only Look Once) và SSD (Single Shot MultiBox Detector). Tuy nhiên, việc triển khai Object Detection trên soft-core RISC-V đòi hỏi tài nguyên tính toán rất lớn.

\subsubsection{Mô hình Phân đoạn (Segmentation)}
Phân loại từng pixel trong ảnh, giúp tách nền và đối tượng một cách chi tiết. Dạng mô hình này (ví dụ: U-Net) thường yêu cầu bộ nhớ lớn để lưu trữ các bản đồ đặc trưng (feature maps) độ phân giải cao, do đó ít được ưu tiên trên các hệ thống FPGA tài nguyên thấp nếu không có bộ tăng tốc phần cứng chuyên dụng.

\subsection{Training và Inference}
Để triển khai AI trên phần cứng, cần phân biệt rõ hai giai đoạn:
\begin{itemize}
    \item \textbf{Huấn luyện (Training):} Quá trình tìm ra bộ trọng số (weights) tối ưu cho mô hình thông qua việc học từ dữ liệu lớn. Quá trình này yêu cầu tính toán dấu phẩy động (Float32/Float16) và thường thực hiện trên Server hoặc GPU mạnh.
    \item \textbf{Suy luận (Inference):} Quá trình sử dụng mô hình đã huấn luyện để đưa ra dự đoán trên dữ liệu mới. Tại thiết bị biên (Edge), suy luận yêu cầu độ trễ thấp và thường sử dụng kỹ thuật lượng tử hóa để chạy trên phần cứng số nguyên (Integer hardware) như RISC-V soft-core.
\end{itemize}

\subsection{Phân loại AI và Mối liên hệ với Deep Learning}
Trí tuệ nhân tạo là một thuật ngữ rộng, bao hàm nhiều cấp độ:
\begin{itemize}
    \item \textbf{AI hẹp (Narrow AI):} Các hệ thống được thiết kế cho một tác vụ cụ thể (như nhận dạng ảnh trong đề tài này).
    \item \textbf{Học máy (Machine Learning - ML):} Tập con của AI, nơi máy tính học từ dữ liệu thay vì được lập trình quy tắc cứng.
    \item \textbf{Học sâu (Deep Learning - DL):} Tập con của ML, sử dụng các mạng nơ-ron nhiều lớp để trích xuất đặc trưng tự động.
\end{itemize}
\textit{Chi tiết về cấu trúc và nguyên lý hoạt động của Mạng nơ-ron nhân tạo (ANN) và Học sâu sẽ được trình bày chuyên sâu tại Mục \ref{sec:neural_network}.}

\subsection{Edge AI}

\subsubsection{Khái niệm}
Edge AI là sự kết hợp giữa thuật toán AI và tính toán biên (Edge Computing), cho phép xử lý dữ liệu cục bộ ngay trên thiết bị phần cứng mà không cần gửi dữ liệu về trung tâm dữ liệu (Cloud).

\subsubsection{Ưu điểm so với Cloud AI}
\begin{itemize}
    \item \textbf{Độ trễ thấp (Low Latency):} Phản hồi tức thời, phù hợp cho các ứng dụng thời gian thực (ví dụ: xe tự lái, robot).
    \item \textbf{Bảo mật và Quyền riêng tư:} Dữ liệu hình ảnh nhạy cảm không rời khỏi thiết bị, giảm nguy cơ bị tấn công đường truyền.
    \item \textbf{Tiết kiệm băng thông:} Không cần truyền tải luồng video liên tục lên máy chủ, giảm tải cho hạ tầng mạng.
    \item \textbf{Độ tin cậy:} Hệ thống vẫn hoạt động bình thường ngay cả khi mất kết nối Internet.
\end{itemize}

\subsubsection{Các Framework phổ biến cho Edge AI}
\begin{itemize}
    \item \textbf{TensorFlow Lite (TFLite):} Phiên bản rút gọn của TensorFlow, tối ưu cho thiết bị di động và nhúng. TFLite hỗ trợ chuyển đổi mô hình sang định dạng FlatBuffer nhẹ và cung cấp các toán tử tối ưu cho vi xử lý ARM và RISC-V.
    \item \textbf{Vitis AI:} Nền tảng phát triển của AMD-Xilinx, chuyên dụng để tối ưu hóa mô hình AI chạy trên các kiến trúc DPU (Deep Learning Processor Unit) của FPGA Xilinx.
    \item \textbf{ONNX Runtime:} ONNX (Open Neural Network Exchange) là định dạng trung gian giúp chuyển đổi mô hình giữa các framework (PyTorch sang TensorFlow). ONNX Runtime cho phép chạy suy luận hiệu năng cao trên nhiều nền tảng phần cứng khác nhau.
\end{itemize}

%=================================================
%=================================================
\section{Mạng Nơ-ron Nhân tạo (Artificial Neural Networks)}
\label{sec:neural_network}

Mạng nơ-ron nhân tạo (ANN) là nền tảng của Deep Learning, 
lấy cảm hứng từ cơ chế truyền dẫn tín hiệu điện hóa giữa các nơ-ron 
trong bộ não sinh học. Trước khi đi sâu vào kiến trúc CNN được sử dụng trong đề tài, 
cần xem xét tổng quan các kiến trúc mạng nơ-ron phổ biến để thấy rõ sự phù hợp 
của từng loại đối với các tác vụ cụ thể.\cite{CNN}

\subsection{Các kiến trúc Mạng Nơ-ron phổ biến}

\subsubsection{Mạng Nơ-ron Truyền thẳng (Feedforward Neural Network - MLP)}
Đây là dạng đơn giản nhất của ANN, còn được gọi là Multi-Layer Perceptron (MLP). Trong kiến trúc này, thông tin di chuyển theo một chiều từ lớp đầu vào (Input), qua các lớp ẩn (Hidden layers) đến lớp đầu ra (Output) mà không có chu trình quay lại.
\begin{itemize}
    \item \textbf{Đặc điểm:} Các nơ-ron ở lớp trước kết nối đầy đủ (fully connected) với nơ-ron lớp sau.
    \item \textbf{Hạn chế với hình ảnh:} Khi xử lý ảnh, MLP làm mất thông tin không gian (spatial structure) do phải duỗi phẳng ảnh thành vector 1 chiều. Ngoài ra, số lượng tham số bùng nổ quá nhanh dẫn đến quá tải bộ nhớ và hiện tượng quá khớp (overfitting).
\end{itemize}

\subsubsection{Mạng Nơ-ron Hồi quy (Recurrent Neural Network - RNN)}
Khác với MLP, RNN có các kết nối quay lại (feedback loops), cho phép thông tin từ các bước trước đó ảnh hưởng đến đầu ra hiện tại. Các biến thể nổi tiếng bao gồm LSTM (Long Short-Term Memory) và GRU.
\begin{itemize}
    \item \textbf{Ứng dụng:} Chuyên biệt cho dữ liệu chuỗi thời gian (Time-series) như giọng nói, văn bản hoặc chứng khoán.
    \item \textbf{Lý do không lựa chọn:} Đối với bài toán nhận dạng vật thể tĩnh trong ảnh, RNN không hiệu quả trong việc trích xuất đặc trưng không gian 2D và chi phí tính toán tuần tự rất khó song song hóa trên FPGA.
\end{itemize}

\subsubsection{Vision Transformers (ViT)}
Đây là kiến trúc tiên tiến nhất hiện nay, áp dụng cơ chế "Sự chú ý" (Self-Attention) từ xử lý ngôn ngữ tự nhiên sang thị giác máy tính. ViT chia ảnh thành các mảnh (patches) và xử lý mối quan hệ toàn cục giữa chúng.
\begin{itemize}
    \item \textbf{Ưu điểm:} Hiệu năng thường vượt trội hơn CNN trên các tập dữ liệu khổng lồ.
    \item \textbf{Thách thức trên Edge AI:} Cơ chế Attention có độ phức tạp tính toán là $O(N^2)$, tiêu tốn tài nguyên phần cứng và bộ nhớ rất lớn, chưa phù hợp để triển khai trên các soft-core RISC-V có xung nhịp thấp như trong phạm vi đề tài này.
\end{itemize}
Và mạng Nơ-ron tích chập (CNN) sẽ được ứng dụng trong đồ án này, sẽ được đề cập ngay sau đây.

\subsection{Mạng Nơ-ron Tích chập (CNN)}
Dựa trên phân tích trên, CNN là kiến trúc cân bằng nhất giữa hiệu năng và chi phí tính toán cho các thiết bị biên. CNN là kiến trúc mạng chuyên biệt cho dữ liệu dạng lưới (grid-like topology) như hình ảnh. Một mạng CNN điển hình bao gồm ba loại lớp chính:

\subsubsection{Lớp Tích chập (Convolutional Layer)}

Đây là lớp quan trọng nhất, thực hiện trích xuất đặc trưng cục bộ và bảo toàn tính chất không gian của ảnh. Phép tính tích chập là việc trượt một cửa sổ nhỏ (Kernel/Filter) trên toàn bộ ảnh đầu vào và thực hiện phép nhân chập.
Về mặt toán học, giá trị đầu ra $Y$ tại vị trí $(i, j)$ được tính như sau:
\begin{equation}
Y[i, j] = \sum_{m} \sum_{n} X[i+m, j+n] \cdot K[m, n] + b
\end{equation}
Trong đó: $X$ là ma trận đầu vào, $K$ là ma trận Kernel, và $b$ là bias. Trên các vi xử lý RISC-V hỗ trợ DSP (như Pulpino), phép toán này có thể được tăng tốc bằng các lệnh SIMD (Single Instruction Multiple Data), cho phép thực hiện 4 phép nhân 8-bit trong một chu kỳ xung nhịp.

\subsubsection{Hàm kích hoạt (Activation Function)}
Hàm kích hoạt giúp đưa tính phi tuyến vào mô hình, cho phép mạng học được các đặc trưng phức tạp. Hàm phổ biến nhất trong CNN là \textbf{ReLU (Rectified Linear Unit)}:
\begin{equation}
f(x) = \max(0, x)
\end{equation}
ReLU có ưu điểm lớn về mặt tính toán trên phần cứng nhúng vì chỉ cần một lệnh so sánh đơn giản (nếu $x < 0$ thì gán bằng 0), không tốn kém tài nguyên như hàm Sigmoid hay Tanh (vốn yêu cầu tính toán số mũ thực rất nặng nề cho soft-core).

\subsubsection{Lớp Gộp (Pooling Layer)}
Dùng để giảm kích thước không gian của bản đồ đặc trưng (down-sampling), giúp giảm số lượng tham số và tính toán, đồng thời giúp mô hình bất biến với các dịch chuyển nhỏ (Translation Invariance) của đối tượng. Hai kỹ thuật phổ biến là Max Pooling (lấy giá trị lớn nhất) và Average Pooling (lấy giá trị trung bình).

\subsection{Kỹ thuật Lượng tử hóa (Quantization)}
Lượng tử hóa là kỹ thuật cốt lõi để triển khai AI trên các vi xử lý RISC-V soft-core không có bộ xử lý dấu phẩy động (FPU) mạnh mẽ.

\subsubsection{Khái niệm}
Kỹ thuật này chuyển đổi các trọng số (weights) và giá trị kích hoạt (activations) từ dạng số thực 32-bit (Float32) sang số nguyên 8-bit (Int8). Công thức chuyển đổi tổng quát:
\begin{equation}
r = S \times (q - Z)
\end{equation}
Trong đó: $r$ là giá trị thực, $q$ là giá trị lượng tử hóa (int8), $S$ là tỉ lệ (Scale factor), và $Z$ là điểm không (Zero-point).

\subsubsection{Lợi ích trên FPGA/RISC-V}
\begin{itemize}
    \item \textbf{Giảm bộ nhớ:} Kích thước mô hình giảm 4 lần (từ 32-bit xuống 8-bit), giúp mô hình vừa vặn với bộ nhớ on-chip (BRAM/TCM) hạn hẹp của FPGA, giảm độ trễ truy cập bộ nhớ ngoài (DDR).
    \item \textbf{Tăng tốc độ:} Các phép tính số học trên số nguyên (Integer Arithmetic) nhanh hơn và tiêu thụ ít năng lượng hơn nhiều so với số thực. Các nhân RISC-V như RI5CY hỗ trợ thực hiện 4 phép nhân cộng Int8 trong một chu kỳ máy.
\end{itemize}

\subsection{Các kỹ thuật tối ưu hóa khác}

\subsubsection{Kỹ thuật Cắt tỉa }
Pruning loại bỏ các kết nối (trọng số) có giá trị gần bằng 0 hoặc ít ảnh hưởng đến kết quả đầu ra. Kỹ thuật này biến ma trận trọng số dày đặc (Dense) thành ma trận thưa (Sparse), giúp giảm đáng kể khối lượng tính toán và bộ nhớ lưu trữ mà không làm giảm nhiều độ chính xác.

\subsubsection{Kỹ thuật im2col (Image-to-Column)}

Để tối ưu hóa phép tính tích chập trên phần cứng, kỹ thuật \textit{im2col} thường được sử dụng. Nó biến đổi ma trận ảnh đầu vào thành các cột, chuyển bài toán tích chập phức tạp (Convolution) thành bài toán nhân ma trận (GEMM - General Matrix Multiply). GEMM là phép toán đã được tối ưu hóa rất tốt trên các thư viện tính toán (như BLAS) và tận dụng triệt để khả năng xử lý song song của phần cứng RISC-V DSP extensions.

\subsection{Các kỹ thuật Lượng tử hóa mở rộng (Extreme Quantization)}
Ngoài kỹ thuật lượng tử hóa Int8 tiêu chuẩn (đã đề cập ở mục 2.5.2), các nghiên cứu hiện đại trên FPGA đang hướng tới việc giảm độ chính xác xuống mức thấp nhất có thể (1-bit hoặc 2-bit) để tối đa hóa hiệu năng phần cứng.

\subsubsection{Mạng Nơ-ron Nhị phân (Binary Neural Networks - BNN)}
BNN là dạng cực đoan nhất của lượng tử hóa, trong đó trọng số (weights) và các giá trị kích hoạt (activations) bị ép buộc chỉ nhận hai giá trị duy nhất: $+1$ và $-1$ (tương ứng với 1 bit).

\textbf{Nguyên lý toán học:}
Quá trình nhị phân hóa thường sử dụng hàm dấu (Sign function):
\begin{equation}
x_b = Sign(x) = \begin{cases}
+1 & \text{nếu } x \ge 0 \\
-1 & \text{nếu } x < 0
\end{cases}
\end{equation}

\textbf{Ưu điểm vượt trội trên FPGA:}
Lợi thế lớn nhất của BNN nằm ở việc thay đổi bản chất phép tính Tích chập:
\begin{itemize}
    \item Trong các mạng thông thường (Float32 hoặc Int8), phép tích chập dựa trên phép nhân và cộng (Multiply-Accumulate - MAC). Điều này tiêu tốn các khối DSP (Digital Signal Processing slices) trên FPGA.
    \item Trong BNN, vì các giá trị chỉ là $\pm 1$, phép nhân số học được thay thế bằng phép toán logic XNOR, và phép cộng được thay thế bằng phép đếm bit (Popcount).
\end{itemize}
\begin{equation}
x \cdot w \approx \text{popcount}(\text{XNOR}(x, w))
\end{equation}
Trên nền tảng Kria KV260, phép toán XNOR-Popcount có thể được thực hiện cực nhanh bằng các bảng tra LUTs (Look-Up Tables) thay vì phải dùng tài nguyên DSP quý giá, giúp tăng thông lượng (Throughput) lên hàng chục lần.

\subsubsection{Mạng Nơ-ron Tam phân (Ternary Weight Networks - TWN)}
Mặc dù BNN rất nhanh, nhưng việc giảm xuống 1-bit thường gây giảm độ chính xác đáng kể (Accuracy drop). TWN là giải pháp lai, cho phép trọng số nhận 3 giá trị: $\{-1, 0, +1\}$.
\begin{itemize}
    \item Giá trị $0$ đóng vai trò quan trọng tạo ra tính thưa (Sparsity). Khi trọng số bằng 0, phép nhân tương ứng có thể được bỏ qua hoàn toàn (Skip operation).
    \item Điều này giúp giảm khối lượng tính toán mà vẫn giữ được độ chính xác gần tiệm cận với mạng Int8.
\end{itemize}

\subsection{Kỹ thuật chưng cất(Knowledge Distillation)}
Để triển khai các mô hình AI nhỏ gọn (như MobileNet, SqueezeNet) lên thiết bị biên mà vẫn đảm bảo độ thông minh, kỹ thuật Chưng cất thường được áp dụng song song với lượng tử hóa.

\textbf{Nguyên lý:}
Mô hình hoạt động dựa trên cơ chế "Thầy - Trò" (Teacher - Student):
\begin{itemize}
    \item \textbf{Teacher Network:} Là một mô hình lớn, phức tạp (ví dụ: ResNet-50), có độ chính xác cao nhưng nặng nề, được huấn luyện trước trên Server.
    \item \textbf{Student Network:} Là mô hình nhỏ gọn (ví dụ: MobileNetV1), là mô hình đích sẽ chạy trên RISC-V/FPGA.
\end{itemize}
Thay vì huấn luyện Student Network trực tiếp trên dữ liệu gốc (Hard labels), ta huấn luyện nó để bắt chước hành vi và xác suất đầu ra (Soft labels) của Teacher Network. Kết quả là Student Network học được các đặc trưng tổng quát tốt hơn và đạt độ chính xác cao hơn so với việc tự học, trong khi kích thước vẫn đủ nhỏ để chạy trên Edge.

\subsubsection{Cơ chế tính toán phần cứng: MAC so với XNOR-Popcount}

Sự khác biệt cốt lõi giữa mạng nơ-ron truyền thống và mạng nhị phân (BNN) nằm ở phép toán cơ bản tại các nơ-ron. Điều này quyết định trực tiếp đến hiệu quả sử dụng tài nguyên trên FPGA.
\begin{figure}[h]
    \centering
    \begin{tikzpicture}[
        node distance=2cm,
        block/.style={rectangle, draw, minimum width=2.5cm, minimum height=1cm, align=center},
        arrow/.style={thick, ->, >=stealth},
        % Tạo thêm style cho Input/Weight để đỡ phải gõ lại nhiều lần
        io_node/.style={align=center} 
    ]
    
    % --- Phần 1: Standard CNN (MAC) ---
    \node (input1) [io_node] {Input ($X$) \\ \textit{32-bit/8-bit}};
    \node (weight1) [right of=input1, xshift=2cm, io_node] {Weight ($W$) \\ \textit{32-bit/8-bit}};
    
    \node (mult) [below of=input1, xshift=2cm] [block] {Multiplier \\ ($\times$)};
    \node (acc) [below of=mult] [block] {Accumulator \\ ($\sum$)};
    \node (dsp) [right of=acc, xshift=2.5cm, text width=3cm, align=left] {\textbf{Tiêu tốn DSP} \\ (Tài nguyên đắt đỏ)};

    \draw[arrow] (input1) |- (mult);
    \draw[arrow] (weight1) |- (mult);
    \draw[arrow] (mult) -- (acc);
    \draw[dashed] (acc) -- (dsp);

    % --- Đường phân cách ---
    % Tinh chỉnh tọa độ y để nằm giữa 2 phần đẹp hơn
    \draw[thick, dashed] (-2, -5.5) -- (8, -5.5);
    \node at (3, -5.8) {\textbf{So sánh với}};

    % --- Phần 2: BNN (XNOR-Popcount) ---
    % Dùng yshift lớn hơn một chút để tránh dính vào đường phân cách
    \node (input2) [below of=acc, yshift=-3.5cm, xshift=-2cm, io_node] {Input ($X_b$) \\ \textit{1-bit}};
    \node (weight2) [right of=input2, xshift=2cm, io_node] {Weight ($W_b$) \\ \textit{1-bit}};
    
    \node (xnor) [below of=input2, xshift=2cm] [block] {XNOR Gate \\ ($\oplus$)};
    \node (pop) [below of=xnor] [block] {Popcount \\ (Đếm bit 1)};
    \node (lut) [right of=pop, xshift=2.5cm, text width=3cm, align=left] {\textbf{Sử dụng LUTs} \\ (Tài nguyên sẵn có)};

    \draw[arrow] (input2) |- (xnor);
    \draw[arrow] (weight2) |- (xnor);
    \draw[arrow] (xnor) -- (pop);
    \draw[dashed] (pop) -- (lut);

    \end{tikzpicture}
    \caption{So sánh luồng dữ liệu giữa phép nhân tích lũy (MAC) truyền thống và phép XNOR-Popcount trong BNN.}
    \label{fig:mac_vs_xnor}
\end{figure}
\paragraph{Phép nhân tích lũy (MAC - Multiply Accumulate):}
Trong các mạng CNN thông thường (như MobileNetV1 sử dụng Int8), phép tính chủ đạo là:
\begin{equation}
    S = \sum_{i=0}^{N-1} x_i \times w_i
\end{equation}
Để thực hiện phép tính này trên FPGA, hệ thống phải sử dụng các khối phần cứng chuyên dụng gọi là DSP Slices. Số lượng DSP Slices trên một chip FPGA là hữu hạn (Kria KV260 có 1,248 DSPs). Khi mô hình quá lớn, số lượng DSP không đủ sẽ trở thành điểm nghẽn (bottleneck), giới hạn khả năng song song hóa.

\paragraph{Phép XNOR-Popcount:}
Trong BNN, do $x_i, w_i \in \{+1, -1\}$, phép nhân được thay thế bằng phép logic đồng nhất (XNOR):
\begin{itemize}
    \item Nếu hai bit giống nhau ($1 \times 1$ hoặc $-1 \times -1$), kết quả là $1$.
    \item Nếu hai bit khác nhau ($1 \times -1$), kết quả là $-1$.
\end{itemize}
Hàm kích hoạt tích chập trở thành:
\begin{equation}
    S = \text{Popcount}(\text{XNOR}(x, w))
\end{equation}
Trên FPGA, phép XNOR có thể được thực hiện trực tiếp bằng các bảng tra Look-Up Tables - đơn vị cơ bản nhất và dồi dào nhất trên FPGA. Một LUT 6 đầu vào (LUT6) có thể xử lý nhiều phép XNOR cùng lúc.
Điều này mang lại hai lợi thế đột phá:
\begin{enumerate}
    \item \textbf{Giải phóng DSP:} Không cần dùng DSP, dành tài nguyên này cho các tác vụ khác.
    \item \textbf{Tối đa hóa song song:} Vì số lượng LUTs lớn hơn rất nhiều so với DSP, ta có thể thực hiện hàng nghìn phép tính đồng thời, tăng tốc độ xử lý (Throughput) lên gấp nhiều lần so với mạng Int8.
\end{enumerate}

